\chapter{Code Snippets}

This chapter presents key code excerpts from the system implementation with explanations.

\section{User Model with Password Hashing (models/User.js)}
\begin{lstlisting}[language=Java, caption={User Schema with bcrypt Pre-save Hook}]
const userSchema = new mongoose.Schema({
  name: { type: String, required: true, trim: true },
  email: { type: String, required: true, unique: true, lowercase: true },
  password: { type: String, required: true, minlength: 6 },
  role: { type: String, enum: ['patient', 'doctor', 'admin'], default: 'patient' },
  isActive: { type: Boolean, default: true }
}, { timestamps: true });

// Hash password before saving
userSchema.pre('save', async function(next) {
  if (!this.isModified('password')) return next();
  const salt = await bcrypt.genSalt(12);
  this.password = await bcrypt.hash(this.password, salt);
  next();
});

// Compare password method
userSchema.methods.comparePassword = async function(candidatePassword) {
  return await bcrypt.compare(candidatePassword, this.password);
};
\end{lstlisting}

\section{JWT Authentication Middleware (middleware/auth.js)}
\begin{lstlisting}[language=Java, caption={JWT Verification and RBAC Middleware}]
const auth = async (req, res, next) => {
  const authHeader = req.header('Authorization');
  if (!authHeader || !authHeader.startsWith('Bearer ')) {
    return res.status(401).json({ error: 'Access denied.' });
  }
  const token = authHeader.replace('Bearer ', '');
  const decoded = jwt.verify(token, process.env.JWT_SECRET);
  const user = await User.findById(decoded.userId).select('-password');
  if (!user || !user.isActive) {
    return res.status(401).json({ error: 'Unauthorized.' });
  }
  req.user = user;
  next();
};

// Role-based access control
const authorize = (...roles) => {
  return (req, res, next) => {
    if (!roles.includes(req.user.role)) {
      return res.status(403).json({ error: 'Access denied.' });
    }
    next();
  };
};
\end{lstlisting}

\section{Alert Threshold Checker (utils/alertChecker.js)}
\begin{lstlisting}[language=Java, caption={Dual-Level Health Threshold Evaluation}]
const THRESHOLDS = {
  heartRate: {
    critical: { min: 40, max: 150 },
    warning:  { min: 50, max: 120 }
  },
  spO2: {
    critical: { min: 0, max: 90 },
    warning:  { min: 0, max: 94 }
  },
  temperature: {
    critical: { min: 35.0, max: 39.5 },
    warning:  { min: 36.0, max: 38.0 }
  }
};

async function checkThresholds(data, io) {
  const { patientId, heartRate, spO2, temperature } = data;
  const alerts = [];
  // Check each vital against thresholds
  if (heartRate < THRESHOLDS.heartRate.critical.min ||
      heartRate > THRESHOLDS.heartRate.critical.max) {
    alerts.push({ patientId, type: 'critical', metric: 'heartRate',
      message: `CRITICAL: Heart rate is ${heartRate} bpm` });
  }
  // Save alerts and broadcast via Socket.IO
  if (alerts.length > 0) {
    const saved = await Alert.insertMany(alerts);
    saved.forEach(a => io.emit('new-alert', a));
  }
  return { status: alerts.length ? 'critical' : 'normal', alerts };
}
\end{lstlisting}

\section{IoT Data Ingestion Endpoint (routes/healthData.js)}
\begin{lstlisting}[language=Java, caption={Real-time IoT Data Processing and Broadcast}]
router.post('/', async (req, res) => {
  const { patientId, deviceId, heartRate, spO2, temperature } = req.body;
  const patient = await Patient.findOne({ patientId });
  const io = req.app.get('io');

  // Check thresholds and generate alerts
  const { status } = await checkThresholds(
    { patientId, patientName: patient?.name, heartRate, spO2, temperature }, io);

  // Save to MongoDB
  const healthData = new HealthData({
    patientId, deviceId, heartRate, spO2, temperature, status });
  await healthData.save();

  // Broadcast via Socket.IO for real-time dashboard
  io.emit('health-data', {
    patientId, heartRate, spO2, temperature, status,
    timestamp: healthData.timestamp });
  res.status(201).json({ data: healthData, status });
});
\end{lstlisting}

\section{IoT Simulator (simulator/iot\_simulator.py)}
\begin{lstlisting}[language=Python, caption={Python Sensor Simulator with Anomaly Injection}]
def simulate_reading(t):
    # Natural sinusoidal variation (circadian rhythm)
    hr_variation = math.sin(t * 0.05) * 5 + random.gauss(0, 2)
    spo2_variation = math.sin(t * 0.03) * 0.8 + random.gauss(0, 0.3)
    temp_variation = math.sin(t * 0.02) * 0.2 + random.gauss(0, 0.05)

    heart_rate = round(base_hr + hr_variation)
    spo2 = round(base_spo2 + spo2_variation, 1)
    temperature = round(base_temp + temp_variation, 1)

    # 8% chance of anomaly per reading
    if random.random() < 0.08:
        anomaly = random.choice(['high_hr','low_hr','low_spo2','high_temp'])
        if anomaly == 'high_hr': heart_rate = random.randint(125, 160)
        elif anomaly == 'low_spo2': spo2 = round(random.uniform(85, 92), 1)
    return heart_rate, spo2, temperature

# Send via HTTP POST every 3 seconds
payload = {"patientId": PATIENT_ID, "heartRate": hr,
           "spO2": spo2, "temperature": temp}
requests.post(SERVER_URL, json=payload, timeout=5)
\end{lstlisting}

\section{React Role-Based Routing (App.jsx)}
\begin{lstlisting}[language=Java, caption={Role-Based Dashboard Routing in React}]
function RoleDashboard() {
  const { user } = useAuth();
  if (user?.role === 'admin') return <AdminDashboard />;
  if (user?.role === 'patient') return <PatientDashboard />;
  return <Dashboard />;  // Doctor dashboard
}

function ProtectedRoute({ children, roles }) {
  const { user } = useAuth();
  if (!user) return <Navigate to="/" replace />;
  if (roles && !roles.includes(user.role))
    return <Navigate to="/dashboard" replace />;
  return children;
}
\end{lstlisting}
